	\begin{center}
		\section{字符串算法}
	\end{center}

	\begin{multicols*}{2}

		\subsection{Manacher 算法}

		Manacher 算法在 $O(n)$ 时间内求出字符串每个位置为中心的最长回文半径。核心思想是利用已经求出的回文区间 $[L, R]$，把当前位置 \texttt{cur} 关于中心对称到 \texttt{R + L - cur}，借助对称位置已经计算好的半径作为初值，避免从 $1$ 开始重新匹配。

		下面分别给出奇回文和偶回文两个版本。

		\begin{minted}{cpp}
string s;
cin >> s;
const int n = s.size();
vector<int> manacher(n);
// 奇回文半径,manacher[cur] 表示以 cur 为中心的最大回文半径
for (int cur = 0, L = 0, R = 0, k = 0; cur < n; cur++)
{
	if (cur <= R)
		k = min(manacher[R + L - cur], R - cur + 1);
	else
		k = 1;
	while (cur + k < n && cur - k >= 0 && s[cur + k] == s[cur - k])
		k++;
	if (cur + k - 1 > R) // 一定保证右端点最右
		R = cur + k - 1, L = cur - k + 1;
	manacher[cur] = k;
}
fill(all(manacher), 0);
// 偶回文半径，manacher[cur] 表示以 {cur - 1, cur} 为中心的最大回文半径
for (int cur = 0, L = 0, R = 0, k = 0; cur < n; cur++)
{
	if (cur <= R)
		k = min(manacher[R + L - cur + 1], R - cur + 1);
	else
		k = 0; // 注意这里要写成 0,因为第一个位置初始不合法
	while (cur + k < n && cur - 1 - k >= 0 && s[cur + k] == s[cur - 1 - k])
		k++;
	if (cur + k - 1 > R) // 一定保证右端点最右
		R = cur + k - 1, L = cur - 1 - k + 1;
	manacher[cur] = k;
}
		\end{minted}

		\subsection{Z 函数（扩展 KMP）}

		Z 函数 $z[i]$ 定义为字符串 $s$ 与其后缀 $s[i..n-1]$ 的最长公共前缀长度，其中约定 $z[0] = n$。算法在 $O(n)$ 时间内维护一个最右匹配区间 $[l, r]$（即所有已计算的 $i$ 中使 $i + z[i] - 1$ 最大的那个 $i$ 给出的区间）：当 \texttt{cur} 落在 $[l, r]$ 内时，可以用对称位置 $\texttt{cur} - l$ 处已经算出的 $z$ 值作为初值，否则从 $0$ 暴力匹配；任何一次扩展都不会让 \texttt{r} 后退，故总匹配次数线性。

		\begin{minted}{cpp}
// Z 函数:求字符串每一个后缀和原字符串的 LCP
auto Z = [&](string s) -> vector<int>
{
	int n = s.size();
	vector<int> Z(n);

	// 默认第一个位置是全匹配
	Z[0] = n;
	for (int cur = 1, l = 0, r = 0, k = 0; cur < n; cur++)
	{
		if (cur <= r)
			k = min(Z[cur - l], r - cur + 1); // 复用维护的匹配区间信息
		else
			k = 0;
		// 暴力扩展
		while (cur + k < n && s[cur + k] == s[k])
			k++;

		// 当且仅当新匹配区间右端点更远的时候更新匹配区间
		if (cur + k - 1 > r)
			l = cur, r = cur + k - 1;
		// 记录 Z 函数的值
		Z[cur] = k;
	}
	return Z;
};
		\end{minted}

		\textbf{函数版}（与本章 \texttt{kmp} 同风格，直接复制即可用；此版把 $z[0]$ 留成 $0$ 而不是 $n$，用到 $z[0]$ 时记得自己特判）：

		\begin{minted}[highlightlines={6}]{cpp}
vector<int> z_function(const string &s) {
	int n = SZ(s);
	vector<int> z(n, 0);
	for (int i = 1, l = 0, r = 0; i < n; ++i) {
		// 落在最右匹配区间内: 用对称位置的答案做初值
		if (i <= r) z[i] = min(z[i - l], r - i + 1);
		while (i + z[i] < n && s[z[i]] == s[z[i] + i]) {
			z[i]++;
		}
		// 只有右端点更远才更新区间
		if (i + z[i] > r) {
			l = i, r = i + z[i] - 1;
		}
	}
	return z;
}
		\end{minted}

		两个版本完全等价：\texttt{min(z[i - l], r - i + 1)} 里，\texttt{z[i - l]} 是对称位置已算好的答案，\texttt{r - i + 1} 是「不能超出已知区间」的截断，取小者当初值后再暴力外扩。

		\textbf{失配字符的用法}：$z[i]$ 匹配到头之后，下标 $z[i]$（前缀侧）与 $i + z[i]$（后缀侧）这两个位置的字符必然\textbf{不同}（除非 $i + z[i] = n$ 已到串尾）。于是每个 $i$ 免费给出一对有序关系 $\big(s[z[i]],\, s[i + z[i]]\big)$，把「后缀 $i$ 与原串的大小/字典序比较」压成一次单字符比较。牛客多校 2026 第 6 场 Full Alphabet 就是用它对每个后缀取出这一对字符，攒成 $26$ 个字母之间的前驱依赖，再交给状压 DP 数拓扑序。

		\begin{minted}{cpp}
auto z = z_function(s);
vector<int> pre(26);              // pre[c]: 字母 c 的前驱依赖集合
for (int i = 0; i < n; ++i) {
	int u = z[i], v = i + z[i];
	if (v >= n) continue;         // 后缀是原串前缀, 无失配位
	u = s[u] - 'a', v = s[v] - 'a';
	if (u == v) continue;
	pre[v] |= 1 << u;             // 这一对的先后由题意决定
}
		\end{minted}

		\subsection{KMP（前缀函数 pi）}

		KMP 的核心是 \textbf{前缀函数}：对串 $A$，$\pi[i]$ 表示 $A[0..i]$ 这一段中最长 \emph{真前缀 = 真后缀} 的长度（约定 $\pi[0] = 0$）。求出 $\pi$ 后做模式匹配的常用技巧是把模式串 $s_2$、文本串 $s_1$ 拼成 \verb|s2 + '#' + s1| 一起跑前缀函数，凡是 $\pi[i] = |s_2|$ 的位置就对应一次匹配。

		\begin{minted}{cpp}
/*
 * 采用 0-based 设计
 * 模式串(s2)+'#'+文本串(s1) 需要这样子传入,如果想要执行 kmp 的话
 * kmp 模板题 AC https://www.luogu.com.cn/record/269142391
 */
vector<ll> kmp(const string &A) {
	// 返回总前缀数组 pi
	// 其定义是 [0, i] 真前后缀相同的最长长度
	// 特别的, pi[0] = 0
	vector<ll> pi(SZ(A));
	for (int i = 1; i < SZ(A); ++i) {
		ll len = pi[i - 1];
		while (len > 0 && A[len] != A[i]) {
			len = pi[len - 1];
		}
		if (A[len] == A[i]) {
			pi[i] = len + 1;
		}
	}
	return pi;
}
		\end{minted}

		典型调用方式：

		\begin{minted}{cpp}
void Solve() {
	string s1, s2;
	cin >> s1 >> s2;
	vector<ll> pi = kmp(s2 + "#" + s1);
	// 输出所有匹配的起始位置
	for (int i = SZ(s2); i < SZ(pi); ++i) {
		if (pi[i] == SZ(s2)) {
			ll ans = i - SZ(s2) - SZ(s2);
			++ans;
			cout << ans << "\n";
		}
	}
	// 输出 s2 第 i 位的最长 border 长度
	for (int i = 0; i < SZ(s2); ++i) {
		cout << pi[i] << " ";
	}
	cout << "\n";
}
		\end{minted}

		\subsection{字符串哈希（mod $2^{61}-1$ 版本）}

		一个用 Mersenne 素数 $2^{61}-1$ 作模的字符串哈希，整数运算用 \verb|__uint128_t| 中转、做一次「低 61 位 + 高位」折叠后只需最多一次减法即可归约到 $[0, M)$。配合 \verb|HS| 类提供的前缀哈希、动态 \verb|push| / \verb|pop|、以及一个简单的 \verb|find| 子串计数接口，覆盖比赛里绝大多数字符串哈希需求。

		\begin{minted}{cpp}
struct U {
	using u128 = __uint128_t;
	static constexpr ll MOD = (1ULL << 61) - 1;

	ll v; // always in [0, MOD)

	U() : v(0) {}
	U(ll x) {
		ll t = x % MOD;
		if (t < 0) t += MOD;
		v = (ll) t;
	}

	static inline ll add(ll a, ll b) {
		ll s = a + b;
		if (s >= MOD) s -= MOD;
		return s;
	}
	static inline ll sub(ll a, ll b) {
		return (a >= b) ? (a - b) : (a + MOD - b);
	}
	static inline ll mul(ll a, ll b) {
		u128 t = (u128) a * b;             // < 2^122
		ll s = ll(t & MOD) + ll(t >> 61);  // 1-step fold
		if (s >= MOD) s -= MOD;            // at most one subtract needed
		return s;
	}

	U &operator+=(const U &o) { v = add(v, o.v); return *this; }
	U &operator-=(const U &o) { v = sub(v, o.v); return *this; }
	U &operator*=(const U &o) { v = mul(v, o.v); return *this; }

	friend U operator+(U a, const U &b) { a += b; return a; }
	friend U operator-(U a, const U &b) { a -= b; return a; }
	friend U operator*(U a, const U &b) { a *= b; return a; }

	friend bool operator==(const U &a, const U &b) { return a.v == b.v; }
	friend bool operator!=(const U &a, const U &b) { return a.v != b.v; }
	friend bool operator< (const U &a, const U &b) { return a.v <  b.v; }
	friend bool operator> (const U &a, const U &b) { return a.v >  b.v; }
	friend bool operator<=(const U &a, const U &b) { return a.v <= b.v; }
	friend bool operator>=(const U &a, const U &b) { return a.v >= b.v; }
};

// base 运行时随机, 防 anti-hash (固定 base 容易被离线预计算碰撞)
const U base = []{
	mt19937_64 rng(chrono::steady_clock::now().time_since_epoch().count());
	uniform_int_distribution<ll> dist(257, U::MOD - 1);
	return U(dist(rng));
}();

// AC https://vjudge.net/solution/65151535 string matching 测试 find
// AC https://www.luogu.com.cn/record/245201000 测试 push / pop / 增强 get
// 字符串哈希约定:越前面是越高位 (和数字一样)
struct HS {
	vector<U> hs, powBase;
	U hashval;
	ll n;
	string s;
	// 内部 0-based, hs[i] 表示 s[0..i] 的哈希
	HS(const string &ls) {
		n = SZ(ls);
		s = ls;
		hs.assign(n, 0);
		powBase.assign(n + 1, 0);
		powBase[0] = 1;
		for (ll i = 1; i <= n; ++i) {
			powBase[i] = powBase[i - 1] * base;
		}
		if (n > 0) {
			hs[0] = U(s[0]);
			for (ll i = 1; i < n; ++i) {
				hs[i] = hs[i - 1] * base + U(s[i]);
			}
			hashval = hs[n - 1];
		} else {
			hashval = 0;
		}
	}

	// get(l, r): 返回 s[l..r] 的哈希 (0-based, inclusive)
	// 自动把越界端点 clamp 到合法范围
	U get(ll l, ll r) {
		if (n == 0) return 0;
		l = max(0ll, l);
		r = min(n - 1, r);
		if (l > r) return 0;
		if (l == 0) return hs[r];
		return hs[r] - hs[l - 1] * powBase[r - l + 1];
	}
	U get(ll l) { return get(l, n - 1); }

	// 返回 ls 在 s 中的出现次数
	ll find(const string &ls) {
		U hxfind = 0;
		ll m = SZ(ls);
		if (m == 0) return n + 1;
		if (m > n) return 0;
		for (ll i = 0; i < m; ++i) {
			hxfind = hxfind * base + ls[i];
		}
		ll res = 0;
		for (ll r = m - 1; r < n; ++r) {
			U hx = get(r - m + 1, r);
			if (hx == hxfind) res++;
		}
		return res;
	}

	void pop() {
		if (n == 0) return;
		n--;
		s.pop_back();
		if (n == 0) {
			hs.clear();
			hashval = 0;
			return;
		}
		hs.pop_back();
		hashval = hs.back();
	}

	void push(char c) {
		s.push_back(c);
		if (n == 0) {
			hs.push_back(U(c));
		} else {
			hs.push_back(hs.back() * base + U(c));
		}
		n++;
		hashval = hs.back();
		if (SZ(powBase) <= n) {
			powBase.push_back(powBase.back() * base);
		}
	}
};
		\end{minted}

		\subsection{子序列自动机}

		子序列自动机预处理 \verb|next[i][c]|：从位置 $i$ 起字符 $c$ 下一次出现的下标（含 $i$）。建出来之后单串子序列匹配可以做到 $O(|t|)$ 一次回答 $t$ 是不是 $s$ 的子序列。下面这一份适配小写字母 / 大写字母 / 数字，换字符集时只需调整 \verb|base| 与第二维大小。

		\begin{minted}{cpp}
struct seqAuto {
	// 子序列自动机 subsequence automaton
	// AC https://www.luogu.com.cn/record/218545124
	vector<vector<int> > next;
	// 从 i 开始, 字符下一次出现的位置 (含 i)
	// 换字符集只需改 base ('a' / 'A' / '0') 和第二维大小
	char base = 'a';

	seqAuto(const string &s) {
		next.assign(SZ(s) + 5, vector<int>(28, -1));
		for (int i = SZ(s) - 1; i >= 0; --i) {
			for (int j = 0; j < 26; ++j) {
				next[i][j] = next[i + 1][j];
			}
			next[i][s[i] - base] = i;
		}
	}

	inline bool isFind(const string &s) {
		ll idx = 0;
		for (int i = 0; i < SZ(s); ++i) {
			if (next[idx][s[i] - base] != -1) {
				idx = next[idx][s[i] - base] + 1;
			} else {
				return false;
			}
		}
		return true;
	}
};
		\end{minted}

		\subsection{后缀自动机（SAM）}

		后缀自动机是接受 $s$ 全部子串的最小 DFA：从初始状态 $0$ 出发沿转移边能走出的字符串，恰好是 $s$ 的子串（走不动 = 不是子串）。状态数 $\le 2n - 1$、转移数 $\le 3n - 4$，在线增量构造，每次 \verb|sam_extend| 均摊 $O(\Sigma)$（这里 $\Sigma = 26$，转移用定长数组存）。

		核心是把「结束位置集合 \verb|endpos| 相同」的子串归成一个状态。同一状态里的串一定互为后缀且长度连续，占据区间 $[\,\verb|len[link]|+1,\ \verb|len|\,]$。后缀链接 \verb|link| 指向「把首字母去掉后掉进另一个 \verb|endpos| 等价类」的那个状态；所有 \verb|link| 反过来看构成一棵以状态 $0$ 为根的树（\verb|endpos| 树 / parent 树），\verb|endpos(v)| 就是 $v$ 子树里所有「前缀状态」的集合。

		\begin{center}
			\begin{forest}
				for tree={draw, rounded corners, font=\small, align=center,
				edge={-Stealth}, l sep=6mm, s sep=3mm, inner sep=3pt}
				[$\varepsilon$\\{\scriptsize len 0}
					[\texttt{a}\\{\scriptsize len 1}]
					[\texttt{b}\\{\scriptsize len 1}
						[\texttt{ab}\\{\scriptsize len 2}]
						[\texttt{abcb}\\{\scriptsize len 4}]
					]
					[\texttt{bc}\\{\scriptsize len 2（clone）}
						[\texttt{abc}\\{\scriptsize len 3}]
						[\texttt{abcbc}\\{\scriptsize len 5}]
					]
				]
			\end{forest}
		\end{center}

		上图是 $s = \texttt{abcbc}$ 的 \verb|link| 树（结点标的是该状态的\emph{最长}串）。读法：结点 \texttt{bc} 表示状态里装着 $\{\texttt{c},\ \texttt{bc}\}$（长度区间 $[\,0+1,\ 2\,]$），它们的 \verb|endpos| 都是 $\{3, 5\}$；而 \texttt{abc} 的 \verb|endpos| 只有 $\{3\}$，所以挂在 \texttt{bc} 下面。这个 \texttt{bc} 正是构造 \texttt{abcbc} 最后一个 \texttt{c} 时分裂出来的 clone。

		\verb|sam_extend| 里唯一麻烦的就是这个 clone：沿 \verb|link| 链往上找到第一个已有 $c$ 转移的 $p$，记 \verb|to = st[p].next[c]|。若 \verb|st[p].len + 1 == st[to].len|，说明 $\mathrm{maxstr}(p) + c$ 就是 \verb|to| 的最长串，\verb|to| 整个都是新前缀的后缀，直接 \verb|link[cur] = to|；否则 \verb|to| 里只有长度 $\le \verb|len[p]|+1$ 的那一半多出了新的结束位置，必须复制出 \verb|clone|（继承 \verb|to| 的全部转移与 \verb|link|，\verb|len| 取 \verb|len[p] + 1|），再把 \verb|link| 链上仍指向 \verb|to| 的那些 $c$ 转移改挂到 \verb|clone|。

		常见套路：不同子串个数 $= \sum_{v \ne 0} \big(\verb|len[v]| - \verb|len[link[v]]|\big)$；某子串的出现次数 = 它所在状态在 \verb|link| 树上的子树里「前缀状态」个数（把状态按 \verb|len| 桶排后倒序把 \verb|cnt| 累加给父亲即可）；两串最长公共子串 = 拿另一串在自动机上走、走不动就沿 \verb|link| 回退并把当前长度压到 \verb|len[link]|。

		两个坑：\textbf{一}是 \verb|st| 要开 \verb|MAXLEN * 2 + 5|（状态数上界 $2n-1$，\verb|MAXLEN| 取单串最长长度）；\textbf{二}是多测每组都要重新 \verb|sam_init()|，且 \verb|sam_extend| 结尾的 \verb|last = cur| 千万别漏——漏了整台自动机就是错的，而且小数据照样跑得出「看着对」的结果。

		\begin{minted}{cpp}
template<int MAXLEN>
struct SufAutomaton {
	struct State {
		int len; // 该状态包含的最长子串的长度
		int link; // 后缀链接，指向包含当前状态最短子串的最长真后缀的状态
		// （把最短子串的第一个字母去掉，也就是所谓的最长真后缀）
		int next[26]; // 字符转移边，假设字符集为小写字母 a-z
	};

	State st[MAXLEN * 2 + 5]; // 状态数上界 2n-1，多开几个保险
	int sz; // 当前自动机中的状态总数，合法编号 0 .. sz-1
	int last; // 指向当前整个字符串完整前缀所在的状态
	static constexpr char BASE = 'a';

	void sam_init() {
		st[0].len = 0;
		st[0].link = -1;
		for (int i = 0; i < 26; ++i) {
			st[0].next[i] = -1;
		}
		sz = 1;
		last = 0;
	}

	void sam_extend(char c) {
		int cur = sz++;
		st[cur].len = st[last].len + 1;
		st[cur].link = -1;
		for (int i = 0; i < 26; ++i) {
			st[cur].next[i] = -1;
		}
		int p = last;
		while (p != -1 && st[p].next[c - BASE] == -1) {
			st[p].next[c - BASE] = cur;
			p = st[p].link;
		}
		if (p == -1) {
			st[cur].link = 0;
		} else {
			int to = st[p].next[c - BASE]; // 注意是 st[p] 不是 st[cur]
			// to 可以作为我们的 link 啊
			// 但是我们需要其为最小串的这个去掉首字母以后的串啊
			if (st[p].len + 1 == st[to].len) {
				// （p 是字符串）p + ‘c’ 就是 to 中的最长子串，因此不需要对这个 to 进行这个分裂
				// 根据 end pos 性质，to 中更短的，也是这个 cur 的这个后缀啊
				st[cur].link = to;
			} else {
				int clone = sz++;
				st[clone].link = st[to].link;
				st[clone].len = st[p].len + 1;
				for (int i = 0; i < 26; ++i) {
					st[clone].next[i] = st[to].next[i];
				}
				while (p != -1 && st[p].next[c - BASE] == to) {
					// 原先的，转移到 to 已经不合适了
					// maxstr(p) + c 的这个 endpos 包含 cur
					// 但是原来的 to 是不包含 cur 的，因此需要更新
					// 更新为这个这个 clone，因为 clone endpos 包含这个 cur
					st[p].next[c - BASE] = clone;
					p = st[p].link;
				}
				st[cur].link = clone;
				// 使用连续性定理（后缀自动机的性质里面提到的），我们知道 to 中包含 len(maxstr(p))+2 的这个字符串啊
				// 因此也接到这个 clone 上面啊
				st[to].link = clone;
			}
		}
		last = cur; // 千万别漏：不推进 last，整台自动机就建错了
	}
};

// 全局开，别放栈上：MAXLEN=1e6 时 st 有 2e6+ 个 State，约 200 MB 级别
// SufAutomaton<1000006> sam;
// sam.sam_init(); for (char c : s) sam.sam_extend(c);
		\end{minted}

		\subsection{Trie 树}

		普通 26 叉 Trie，节点上挂 \verb|cnt| 表示有多少串经过该节点；\verb|query(s)| 返回已有集合中所有串与 $s$ 的最长公共前缀长度之和（每个串按其与 $s$ 的 LCP 长度计入一次）。\verb|remove| 用懒计数撤回，不释放节点。

		\begin{minted}{cpp}
/*
 * AC https://codeforces.com/gym/105941/submission/367570263
 * 统计已有集合中所有字符串与 s 的相同前缀长度之和
 */
class Trie {
	vector<array<int, 26>> tree;
	vector<int> cnt;
	int idx, n;

public:
	Trie(int N = 2e6 + 1) : n(N), tree(N), cnt(N), idx(0) {}

	void insert(const string s) {
		int cur = 0;
		for (const char ch : s) {
			int pos = ch - 'a';
			if (tree[cur][pos] == 0) {
				tree[cur][pos] = ++idx;
			}
			cur = tree[cur][pos];
			cnt[cur]++;
		}
	}

	void remove(const string s) {
		int cur = 0;
		for (const char ch : s) {
			int pos = ch - 'a';
			if (tree[cur][pos] == 0) {
				assert(0);
			}
			cur = tree[cur][pos];
			cnt[cur]--;
		}
	}

	ll query(const string &s) {
		int cur = 0;
		ll res = 0;
		for (const char ch : s) {
			int pos = ch - 'a';
			if (tree[cur][pos] == 0) break;
			cur = tree[cur][pos];
			res += cnt[cur];
		}
		return res;
	}
};
		\end{minted}

		\subsection{01 Trie（XOR Trie）}

		处理一族 \verb|unsigned long long| 上的最大 / 最小异或问题。\verb|max_xor(x)| / \verb|min_xor(x)| 返回与 $x$ 异或的极值；\verb|argmax_xor| / \verb|argmin_xor| 返回取得这个极值的 \emph{原数} 本身。\verb|erase| 走「沿路径减 \verb|cnt|，遇到 \verb|cnt| 归零的子树则懒摘除指针」的写法，避免正经回收带来的额外管理。

		\begin{minted}{cpp}
// AC https://judge.yosupo.jp/submission/308823 min_xor / erase / insert
// AC https://codeforces.com/contest/706/submission/334623853 argmax / max
// AC https://www.codechef.com/viewsolution/1183578991 综合
struct XORTrie {
	int B; // 最高位 (含)
	struct Node {
		int nxt[2];
		int cnt; // 经过该节点的数的个数 (叶子 = 出现次数)
		Node() { nxt[0] = nxt[1] = -1; cnt = 0; }
	};
	vector<Node> tr;
	explicit XORTrie(int max_bit = 63) : B(max_bit) {
		tr.reserve(1 << 20);
		tr.emplace_back(); // root
	}
	void clear() { tr.clear(); tr.emplace_back(); }

	inline void insert(unsigned long long x) {
		int u = 0;
		for (int b = B; b >= 0; --b) {
			int bit = int((x >> b) & 1ull);
			if (tr[u].nxt[bit] == -1) {
				tr[u].nxt[bit] = (int) tr.size();
				tr.emplace_back();
			}
			u = tr[u].nxt[bit];
			++tr[u].cnt;
		}
	}
	inline int count(unsigned long long x) const {
		int u = 0;
		for (int b = B; b >= 0; --b) {
			int bit = int((x >> b) & 1ull);
			int v = tr[u].nxt[bit];
			if (v == -1) return 0;
			u = v;
		}
		return tr[u].cnt;
	}
	inline bool erase(unsigned long long x) {
		if (tr.size() == 1) return false;
		static int pathBit[70];
		int u = 0, depth = 0;
		for (int b = B; b >= 0; --b) {
			int bit = int((x >> b) & 1ull);
			int v = tr[u].nxt[bit];
			if (v == -1 || tr[v].cnt == 0) return false;
			pathBit[depth++] = bit;
			u = v;
		}
		if (tr[u].cnt == 0) return false;
		u = 0;
		for (int i = 0; i < depth; ++i) {
			int bit = pathBit[i];
			int v = tr[u].nxt[bit];
			--tr[v].cnt;
			if (tr[v].cnt == 0) tr[u].nxt[bit] = -1; // 懒回收
			u = v;
		}
		return true;
	}
	inline unsigned long long max_xor(unsigned long long x) const {
		if (tr.size() == 1) return 0ull;
		int u = 0;
		unsigned long long ans = 0;
		for (int b = B; b >= 0; --b) {
			int bit    = int((x >> b) & 1ull);
			int prefer = tr[u].nxt[bit ^ 1];
			int same   = tr[u].nxt[bit];
			if (prefer != -1 && tr[prefer].cnt > 0) {
				u = prefer;
				ans |= (1ull << b);
			} else {
				u = same; // 至少有一条有效路
			}
		}
		return ans;
	}
	inline unsigned long long min_xor(unsigned long long x) const {
		if (tr.size() == 1) return 0ull;
		int u = 0;
		unsigned long long ans = 0;
		for (int b = B; b >= 0; --b) {
			int bit   = int((x >> b) & 1ull);
			int same  = tr[u].nxt[bit];
			int other = tr[u].nxt[bit ^ 1];
			if (same != -1 && tr[same].cnt > 0) {
				u = same;
			} else {
				u = other;
				ans |= (1ull << b);
			}
		}
		return ans;
	}
	inline unsigned long long argmax_xor(unsigned long long y) const {
		if (tr.size() == 1) return 0ull;
		int u = 0;
		unsigned long long chosen = 0;
		for (int b = B; b >= 0; --b) {
			int bit    = int((y >> b) & 1ull);
			int prefer = tr[u].nxt[bit ^ 1];
			int same   = tr[u].nxt[bit];
			if (prefer != -1 && tr[prefer].cnt > 0) {
				u = prefer;
				chosen |= (unsigned long long) (bit ^ 1) << b;
			} else {
				u = same;
				chosen |= (unsigned long long) bit << b;
			}
		}
		return chosen;
	}
	inline unsigned long long argmin_xor(unsigned long long x) const {
		if (tr.size() == 1) return 0ull;
		int u = 0;
		unsigned long long chosen = 0;
		for (int b = B; b >= 0; --b) {
			int bit   = int((x >> b) & 1ull);
			int same  = tr[u].nxt[bit];
			int other = tr[u].nxt[bit ^ 1];
			if (same != -1 && tr[same].cnt > 0) {
				u = same; // 让该位 XOR=0
				if (bit) chosen |= (1ull << b);
			} else {
				u = other; // 只能走相反位, 该位 XOR=1
				if (bit ^ 1) chosen |= (1ull << b);
			}
		}
		return chosen;
	}
} tr;
// 处理无符号数
		\end{minted}

		\subsection{字符串分词与输出技术（大模拟读入）}

		\verb|getline + stringstream| 是处理一行内含混合分隔符的标准搭配：先 \verb|getline| 整行，再丢进 \verb|stringstream| 按空白 / 自定义字符切分。下面两段一段是以 \verb|':' / '+'| 等多种字符做模板化分词的「跨天时间均值」例题，另一段是只用 \verb|'+'| 切的最简加法器。

		\begin{minted}{cpp}
// 跨天时间区间均值: 行内固定用 ':' / 空格分隔, 末尾可选 (+day)
inline void solve() {
	init();
	FOR(i, 1, 2) {
		string line; getline(cin, line);
		stringstream sst(line); char fen;
		sst >> hh[1] >> fen >> mm[1] >> fen >> ss[1]
		    >> hh[2] >> fen >> mm[2] >> fen >> ss[2];
		string dayst; ll day = 0;
		if (sst >> dayst) {
			stringstream st(dayst);
			st >> fen >> fen >> day;
		}
		sec[i] = day * 86400 + 3600 * (hh[2] - hh[1])
		         + 60 * (mm[2] - mm[1]) + (ss[2] - ss[1]);
	}
	ll fsec = (sec[1] + sec[2]) / 2;
	ll h = fsec / 3600; fsec %= 3600;
	ll m = fsec / 60;   fsec %= 60;
	ll s = fsec;
	cout << setw(2) << setfill('0') << h << ":"
	     << setw(2) << setfill('0') << m << ":"
	     << setw(2) << setfill('0') << s << "\n";
}
		\end{minted}

		\begin{minted}{cpp}
// 用 '+' 作分词符的简单加法器
inline void __() {
	stringstream ss(line);
	string num;
	ll ans = 0;
	while (getline(ss, num, '+')) {
		ans += stoll(num);
	}
	cout << ans << "\n";
}
		\end{minted}

	\end{multicols*}

